\input format.tex

\usepackage{graphicx}
\graphicspath{{cores/}}
\usepackage{setspace}
\usepackage{colortbl,array,booktabs}

\newcolumntype{K}[1]{>{\centering\arraybackslash}p{#1}}

\begin{document}

\vspace*{5mm}

\color{gray2}
\noindent\fontsize{16pt}{\baselineskip}\selectfont{肠道菌群基因检测结果}

%% 各章节
\setlength{\arrayrulewidth}{1pt}
\fontsize{8pt}{11pt}\selectfont

\vspace*{4mm}

\fontsize{10pt}{13pt}\selectfont
\noindent\begin{tabular}{p{4cm}p{6cm}p{4cm}}
\parbox[c]{\hsize}{\vskip5pt 姓名：荔湾 \vskip5pt} & 年龄：35 &
性别：男
\\
\parbox[c]{\hsize}{\vskip5pt 样本编码：GI00001 \vskip5pt} & 送检机构：天津河东民众健检门诊部有限责任公司 & 报告日期：2019-02-19\\
\end{tabular}

\vspace*{4mm}
\begin{spacing}{1.2}
\noindent{\fontsize{12pt}{11pt}\selectfont {主要有益菌检测结果}}
\end{spacing}
\arrayrulecolor{gray2}
\setlength{\arrayrulewidth}{.5pt}

\vspace*{3mm}
\fontsize{9pt}{11pt}\selectfont
\noindent
\begin{tabular}{m{3.4cm}<{\centering}m{2.8cm}<{\centering}m{2.8cm}<{\centering}m{5.2cm}<{\centering}}
  \parbox[c][0.6cm]{\hsize}{\cellcolor{topcolor}\color{white}\centering\bf 主要有益菌属} &
  \parbox[c][0.6cm]{\hsize}{\cellcolor{topcolor}\color{white}\centering\bf 本次检测值} &
  \parbox[c][0.6cm]{\hsize}{\cellcolor{topcolor}\color{white}\centering\bf 参考范围} &
  \parbox[c][0.6cm]{\hsize}{\cellcolor{topcolor}\color{white}\centering\bf 功能} \\
\end{tabular}

\vspace*{-3.5mm}
\fontsize{8.8pt}{11pt}\selectfont
\begin{longtable}{|m{3.4cm}<{\centering}|m{2.8cm}<{\centering}|m{2.8cm}<{\centering}|m{5.1cm}<{\centering}|}
\hline
\parbox[c]{\hsize}{\vskip7pt\centering 乳酸杆菌属 \vskip7pt} & 
0.02
\raisebox{-0.3ex}{\includegraphics[scale=1]{redarrow_down.pdf}}
& ≥ 0.04
& 抑制有害菌的生长，增强人体免疫力 \\
\hline
\parbox[c]{\hsize}{\vskip7pt\centering 双歧杆菌属 \vskip7pt} & 
0.01
\raisebox{-0.3ex}{\includegraphics[scale=1]{redarrow_down.pdf}}
& ≥ 0.21
& 抵抗有害菌入侵，促进微量元素吸收 \\
\hline
\parbox[c]{\hsize}{\vskip7pt\centering 阿克曼氏菌属 \vskip7pt} & 
未检出
\raisebox{-0.3ex}{\includegraphics[scale=1]{redarrow_down.pdf}}
& ≥ 0.01
& 促进脂肪代谢，预防肥胖 \\
\hline
\parbox[c]{\hsize}{\vskip7pt\centering 柔嫩梭菌属 \vskip7pt} & 
13.87
\raisebox{-0.3ex}{\includegraphics[scale=1]{redarrow_down.pdf}}
& ≥ 27.24
& 产丁酸，抑制肠道炎症 \\
\hline
\end{longtable}

\vspace*{-2mm}
\noindent {结果说明：\raisebox{-0.3ex}{\includegraphics[scale=1]{redarrow_down.pdf}} \ 检测值相比参考人群偏低}

\vspace*{2mm}
\noindent {结果分析：
乳酸杆菌属、双歧杆菌属、阿克曼氏菌属、柔嫩梭菌属结果均正常；\\\\
}

%\vspace*{-3mm}
\noindent\fontsize{12pt}{12pt}\selectfont {菌群调理建议} \\

\vspace*{-2mm}
\begin{spacing}{1.3}
\noindent\fontsize{9pt}{11pt}\selectfont
本次结果各主要有益菌含量均正常，以下建议帮助维持并增强有益菌及肠道健康。\\
\end{spacing}

\setlength{\arrayrulewidth}{.5pt}
\vspace*{-2mm}
\begin{center}
\begin{tabular}{K{2.5cm}K{12.3cm}}
\hline

\parbox[c][3cm]{.95\hsize}{
\noindent
\includegraphics[scale=0.8]{baohulifoodwithvline.pdf}
}
 &
%\hspace*{1mm}
\parbox{\hsize}{
\vspace*{0.5mm}
\begin{spacing}{1.5}
\fontsize{10pt}{10pt}\selectfont{\bf {膳食建议 \\}}
{\includegraphics[scale=1.0]{xiaoyuandian.pdf}\fontsize{9pt}{10pt}\selectfont \ 建议保持定时定量进餐的好习惯，早餐不可忽略。食物种类新鲜多样，交替食用。主食搭配适量粗粮，多食用各种新鲜的蔬菜水果，深绿色蔬菜至少占蔬菜总摄入量的一半。}
\end{spacing}
} \\
\hline

\parbox[c][4.5cm]{.95\hsize}{
\noindent
\includegraphics[scale=0.8]{baohuliprobioticswithvline.pdf}
}
 &
%\hspace*{1mm}
\parbox{\hsize}{
\vspace*{0.5mm}
\begin{spacing}{1.5}
\fontsize{10pt}{10pt}\selectfont{\bf {肠道调节建议 \\}}
{\includegraphics[scale=1.0]{xiaoyuandian.pdf}\fontsize{9pt}{10pt}\selectfont \ {\bf {补充益生菌：}} 您平时可适量多喝酸奶。酸奶营养丰富易吸收，含有大量活性益生菌的酸奶（如保加利亚乳杆菌、BB-12），能帮助调节肠道酸碱值，抑制有害菌的生长，增强人体免疫力。 \\}
{\includegraphics[scale=1.0]{xiaoyuandian.pdf}\fontsize{9pt}{10pt}\selectfont \ {\bf {补充益生元：}} 您可选择性补充益生元，如低聚果糖、低聚半乳糖、菊粉等，能帮助益生菌在肠道内增殖，增加有益物质的产生，有利于增强肠道屏障功能，促进肠道健康。}
\end{spacing}
} \\
\hline

\end{tabular}
\end{center}


\end{document}

